%% TL-NN Paper — Transform, Linearize, Neural Network
%% A Hybrid Method for Nonlinear Fredholm Integral Equations on Large Intervals
%% Author: Ilyes SEDKA, Ammar KHELLAF
%% Target: Applied Mathematics and Computation (Q2, Elsevier)

\documentclass[12pt, a4paper]{article}

%% ============================================================
%% PACKAGES
%% ============================================================
\usepackage[top=2.5cm, bottom=2.5cm, left=2.5cm, right=2.5cm]{geometry}
\usepackage{amsmath, amssymb, amsthm}
\usepackage{mathtools}
\usepackage{graphicx}
\usepackage{booktabs}
\usepackage{array}
\usepackage{multirow}
\usepackage{hyperref}
\usepackage{cite}
\usepackage{lineno}
\usepackage{setspace}
\usepackage{xcolor}
\usepackage{algorithm}
\usepackage{algpseudocode}
\usepackage{caption}
\usepackage{subcaption}

%% ============================================================
%% THEOREM ENVIRONMENTS
%% ============================================================
\newtheorem{theorem}{Theorem}[section]
\newtheorem{proposition}[theorem]{Proposition}
\newtheorem{lemma}[theorem]{Lemma}
\newtheorem{corollary}[theorem]{Corollary}
\newtheorem{definition}[theorem]{Definition}
\newtheorem{remark}[theorem]{Remark}
\newtheorem{assumption}[theorem]{Assumption}

%% ============================================================
%% CUSTOM COMMANDS
%% ============================================================
\newcommand{\norm}[1]{\left\| #1 \right\|}
\newcommand{\abs}[1]{\left| #1 \right|}
%\newcommand{\Psi}{\boldsymbol{\Psi}}
%\newcommand{\Phi}{\boldsymbol{\Phi}}
\newcommand{\HNN}{H_{\mathrm{NN}}^{(k)}}
\newcommand{\Hn}{H_{n}^{(k)}}
\newcommand{\DT}{\mathrm{DT}}
\newcommand{\DTNN}{\mathrm{DT}_{\mathrm{NN}}}
\newcommand{\DTn}{\mathrm{DT}_{n}}
\newcommand{\ftheta}{f_{\theta}^{(k,lp)}}
\newcommand{\epsNN}{\varepsilon_{\mathrm{NN}}}
\newcommand{\alphNN}{\alpha_{\mathrm{NN}}}
\newcommand{\varrhoNN}{\varrho_{\mathrm{NN}}}
\newcommand{\R}{\mathbb{R}}
\newcommand{\N}{\mathbb{N}}

%% ============================================================
%% DOCUMENT
%% ============================================================
\begin{document}

\doublespacing

%% ------------------------------------------------------------
%% TITLE PAGE
%% ------------------------------------------------------------
\title{\textbf{A Hybrid Transform--Linearize--Neural Network Method\\
for Nonlinear Fredholm Integral Equations\\
on Large Intervals}}

\author{
  Ilyes SEDKA$^{1,2}$ \and Ammar KHELLAF$^{3}$
}

\date{}
\maketitle

\begin{center}
{\small
$^{1}$ Amar Telidji University, Laghouat, Algeria \\
$^{2}$ Laboratory of Applied Mathematics and Modelling (LMAM),
University 8 Mai 1945 Guelma, Algeria \\
$^{3}$ National Polytechnic College of Constantine, Algeria \\[6pt]
\textit{E-mail:} \href{mailto:ilyes.sedka@lagh-univ.dz}{ilyes.sedka@lagh-univ.dz},
\href{mailto:amarlasix@gmail.com}{amarlasix@gmail.com}
}
\end{center}

\vspace{1cm}

%% ------------------------------------------------------------
%% ABSTRACT
%% ------------------------------------------------------------
\begin{abstract}
We propose a novel hybrid numerical method, called TL-NN
(Transform, Linearize, Neural Network), for solving nonlinear
Fredholm integral equations of the second kind defined on large
intervals $[0, \tau]$ with $\tau \gg 1$.
The method builds directly on the TLD (Transform, Linearize,
Discretize) algorithm introduced in \cite{sedka2025},
preserving its first two phases --- domain decomposition and
Newton--Kantorovich linearization --- while replacing the
Nystr\"{o}m discretization phase with a physics-informed
neural network approximator.
At each Newton iteration, a small feedforward neural network
is trained to approximate the Fr\'{e}chet derivative of the
nonlinear integral operator, yielding matrix entries that
replace those formerly produced by quadrature.
The convergence analysis of TLD carries over to TL-NN with
the Nystr\"{o}m approximation error $\delta_{n,N}$ replaced
by the neural network approximation error $\epsNN$, and we
discuss conditions under which $\epsNN < \delta_{n,N}$ on
large intervals.
Numerical experiments on three benchmark examples from
\cite{sedka2025} for $\tau \in \{20, 50, 80\}$ demonstrate
that TL-NN achieves accuracy comparable to TLD while reducing
the sensitivity to quadrature grid density on large domains.
\end{abstract}

\textbf{Keywords:} Nonlinear Fredholm integral equations;
large interval; Newton--Kantorovich method; neural network
approximation; physics-informed neural networks;
hybrid numerical method; domain decomposition.

\textbf{MSC (2020):} 45G15, 65J15, 68T07, 49M15.

\newpage

%% ============================================================
\section{Introduction}
\label{sec:intro}
%% ============================================================

Nonlinear Fredholm integral equations of the second kind arise
naturally in mathematical models of heat transfer, fluid
dynamics, population dynamics, and electromagnetic scattering
\cite{atkinson1997, ladopoulos2000}.
Their numerical solution is challenged by three simultaneous
difficulties: the nonlinearity of the kernel with respect to
the unknown function, the nonlocal character of the integral
operator whose action at any point depends on the global
behavior of the solution, and --- when the domain is large ---
the high computational cost of maintaining sufficient accuracy
across the entire interval.

Classical numerical approaches, including the
Newton--Kantorovich method combined with Nystr\"{o}m
discretization \cite{atkinson1997, grammont2013, grammont2014},
are well-established and come with rigorous convergence
guarantees.
However, on large intervals their performance degrades
because the number of quadrature nodes required for an
accurate Nystr\"{o}m approximation grows with the domain size,
making the resulting linear systems expensive to solve at
each Newton iteration.

To address this, the TLD (Transform, Linearize, Discretize)
algorithm was introduced in \cite{sedka2025}, where the large
interval $[0, \tau]$ is first decomposed into $N$ smaller
subintervals (the Transform phase), the nonlinear system is
then linearized by the Newton--Kantorovich method (the
Linearize phase), and the resulting linear integral system is
finally discretized by the Nystr\"{o}m method (the Discretize
phase).
The TLD algorithm is supported by a complete convergence
theorem guaranteeing that the approximate solution converges
to the exact solution at rate $(1 + \alpha_n)/2$, regardless
of the size of $\tau$, provided $n$ is large enough.

Independently, the emergence of physics-informed neural
networks (PINNs) \cite{raissi2019} and operator learning
methods \cite{lu2021, kovachki2023} has opened new avenues
for solving integral and differential equations using machine
learning.
These methods parametrize the solution by a neural network
and minimize a residual-based loss, bypassing the need for
explicit quadrature.
While they offer flexibility and generalization capability,
they currently lack the rigorous convergence guarantees that
characterize classical numerical methods \cite{lu2021}.

The present paper proposes a hybrid approach that unites the
theoretical soundness of TLD with the approximation
flexibility of neural networks.
We call the resulting method TL-NN (Transform, Linearize,
Neural Network).
The key idea is simple: keep the first two phases of TLD
exactly as in \cite{sedka2025}, and replace only the
Nystr\"{o}m discretization phase with a small feedforward
neural network trained, at each Newton iteration, to
approximate the Fr\'{e}chet derivative of the integral
operator.
The matrix structure, the linear system, and the iteration
framework remain identical to TLD.
The neural network approximation error $\epsNN$ plays the
same structural role as the Nystr\"{o}m error $\delta_{n,N}$,
and the convergence analysis of \cite{sedka2025} carries over
directly.

The specific advantage of TL-NN over TLD concerns the
behavior on large intervals with complex or oscillatory
kernels.
The Nystr\"{o}m method distributes its approximation capacity
uniformly over a fixed quadrature grid, requiring many nodes
to resolve the kernel derivative accurately everywhere on a
large domain.
A neural network trained on well-chosen collocation points
distributes its capacity adaptively, achieving comparable
approximation quality with fewer evaluation points once
trained.
After training, evaluation of the network at any point costs
a fixed amount independent of $\tau$.

\textbf{Contributions of this paper.}
\begin{enumerate}
  \item We formulate the TL-NN algorithm, a hybrid method
    that integrates a neural network approximator into the
    Newton--Kantorovich iteration framework of TLD for
    nonlinear Fredholm integral equations on large intervals.
  \item We carry over the convergence analysis of
    \cite{sedka2025} to TL-NN, expressing the convergence
    rate in terms of the neural network approximation error
    $\epsNN$ and identifying conditions under which TL-NN
    improves on TLD.
  \item We provide numerical experiments on three benchmark
    examples from \cite{sedka2025}, producing direct
    error table comparisons between TLD and TL-NN for
    $\tau \in \{20, 50, 80\}$ and $n \in \{5, 10, 30, 50, 100\}$.
\end{enumerate}

\textbf{Organization of the paper.}
Section~\ref{sec:preliminaries} recalls the TLD framework
and introduces the PINN approximation approach.
Section~\ref{sec:tlnn} presents the TL-NN algorithm in full
detail.
Section~\ref{sec:convergence} develops the convergence
discussion.
Section~\ref{sec:numerics} reports the numerical experiments.
Section~\ref{sec:conclusion} concludes and discusses future
directions.

%% ============================================================
\section{Preliminaries}
\label{sec:preliminaries}
%% ============================================================

\subsection{The TLD Framework}
\label{subsec:tld}

We recall the main elements of the TLD algorithm from
\cite{sedka2025}, which constitutes the foundation of the
present work.

\subsubsection{The Problem}

Let $\tau \gg 1$ be a large real number.
Consider the nonlinear Fredholm integral equation of the
second kind:
\begin{equation}
  \lambda \psi(t) = \int_0^\tau \kappa(t, s, \psi(s)) \, ds
  + g(t), \quad t \in [0, \tau],
  \label{eq:main}
\end{equation}
where $\psi : [0, \tau] \to \R$ is the unknown function,
$\kappa : [0, \tau]^2 \times \R \to \R$ is a nonlinear kernel
differentiable with respect to its third argument,
$g \in C([0, \tau], \R)$ is a given function, and
$\lambda \in \R^* = \R \setminus \{0\}$.

The Banach space of solutions is:
\begin{equation*}
  \mathcal{X} = \prod_{p=1}^{N} C^1([\tau_{p-1}, \tau_p], \R),
  \quad
  \norm{\Phi}_\mathcal{X} = \sum_{p=1}^{N}
  \norm{\phi_p}_\infty
  = \sum_{p=1}^{N} \max_{t \in [\tau_{p-1}, \tau_p]}
  \abs{\phi_p(t)},
\end{equation*}
equipped with the norm arising from the interval
decomposition defined below.

\subsubsection{The Transform Phase}

For $N \in \N^*$, divide $[0, \tau]$ into $N$ equal
subintervals:
\begin{equation}
  [0, \tau] = \bigcup_{p=1}^{N} [\tau_{p-1}, \tau_p],
  \quad
  \tau_0 = 0, \quad \tau_N = \tau,
  \quad \tau_p = p \cdot \frac{\tau}{N}.
  \label{eq:decomp}
\end{equation}

For each $l = 1, \ldots, N$, define the local functions:
\begin{equation*}
  \psi_l(t) = \psi(t), \quad t \in [\tau_{l-1}, \tau_l],
  \qquad
  g_l(t) = g(t),
  \qquad
  \kappa_l(t, s, \cdot) = \kappa(t, s, \cdot).
\end{equation*}

Equation~\eqref{eq:main} is equivalent to the system of $N$
coupled equations:
\begin{equation}
  \lambda \psi_l(t) = \sum_{p=1}^{N}
  \int_{\tau_{p-1}}^{\tau_p}
  \kappa_l(t, s, \psi_p(s)) \, ds + g_l(t),
  \quad t \in [\tau_{l-1}, \tau_l],
  \quad l = 1, \ldots, N.
  \label{eq:system}
\end{equation}

Defining $\boldsymbol{\Psi} = (\psi_1, \ldots, \psi_N)^\top
\in \mathcal{X}$,
$\mathbf{G} = (g_1, \ldots, g_N)^\top$, and the nonlinear
operator $\mathbf{K} = (K_1, \ldots, K_N)^\top$ by:
\begin{equation}
  K_l(\boldsymbol{\Phi})(t) = \sum_{p=1}^{N}
  \int_{\tau_{p-1}}^{\tau_p}
  \kappa_l(t, s, \phi_p(s)) \, ds,
  \quad t \in [\tau_{l-1}, \tau_l],
  \label{eq:Kl}
\end{equation}
the system~\eqref{eq:system} takes the compact operator form:
\begin{equation}
  \text{Find } \boldsymbol{\Psi} \in \mathcal{X} :
  \quad F(\boldsymbol{\Psi}) =
  \lambda \boldsymbol{\Psi} - \mathbf{K}(\boldsymbol{\Psi})
  - \mathbf{G} = \mathbf{0}_\mathcal{X}.
  \label{eq:operator}
\end{equation}

\subsubsection{Assumptions}

We work under the following assumptions, identical to those
of \cite{sedka2025}.

\begin{assumption}
\label{ass:main}
For all $l, p = 1, \ldots, N$:
\begin{itemize}
  \item[\emph{(i)}] Problem~\eqref{eq:operator} has a unique
    solution $\boldsymbol{\Psi} \in \mathcal{X}$.
  \item[\emph{(ii)}] $(\lambda I - \DT(\boldsymbol{\Psi}))$
    is invertible and there exists $\nu > 0$ such that
    $\norm{(\lambda I - \DT(\boldsymbol{\Psi}))^{-1}} \leq \nu < \infty$.
  \item[\emph{(iii)}] There exists $\gamma_p > 0$ such that
    \begin{equation*}
      \left\|
        \frac{\partial \kappa_l}{\partial \phi_p}(t, s, \phi_p(s))
        -
        \frac{\partial \kappa_l}{\partial \phi_p}(t, s, \varphi_p(s))
      \right\|_\infty
      \leq \gamma_p \norm{\phi_p - \varphi_p}_\infty.
    \end{equation*}
\end{itemize}
\end{assumption}

\subsubsection{The Linearize Phase}
\label{subsubsec:linearize}

The Newton--Kantorovich method linearizes
\eqref{eq:operator} as:
\begin{equation}
  \bigl(\lambda I - \DT(\boldsymbol{\Psi}^{(k)})\bigr)
  \bigl(\boldsymbol{\Psi}^{(k+1)} - \boldsymbol{\Psi}^{(k)}\bigr)
  = -F(\boldsymbol{\Psi}^{(k)}),
  \quad \boldsymbol{\Psi}^{(0)} \in \mathcal{X},
  \quad k = 0, 1, 2, \ldots
  \label{eq:newton}
\end{equation}
where $\DT(\boldsymbol{\Psi}^{(k)})$ is the Fr\'{e}chet
derivative of $\mathbf{K}$ at $\boldsymbol{\Psi}^{(k)}$,
whose $(l,p)$-block component is the linear integral operator:
\begin{equation}
  [T_{lp}(\boldsymbol{\Psi}^{(k)}) \phi_p](t) =
  \int_{\tau_{p-1}}^{\tau_p}
  \frac{\partial \kappa_l}{\partial \psi_p}
  \bigl(t, s, \psi_p^{(k)}(s)\bigr) \cdot \phi_p(s) \, ds,
  \quad t \in [\tau_{l-1}, \tau_l].
  \label{eq:Tlp}
\end{equation}

The right-hand side at iteration $k$ is:
\begin{equation}
  f_l^{(k)}(t) =
  - \sum_{p=1}^{N}
  \int_{\tau_{p-1}}^{\tau_p}
  \frac{\partial \kappa_l}{\partial \psi_p}
  \bigl(t, s, \psi_p^{(k)}(s)\bigr)
  \cdot \psi_p^{(k)}(s) \, ds
  + K_l(\boldsymbol{\Psi}^{(k)})(t) + g_l(t).
  \label{eq:rhs}
\end{equation}

\subsubsection{The Discretize Phase (Nystr\"{o}m)}

For each $l = 1, \ldots, N$, consider the equidistant
subdivision $\Omega_n^l$ of $[\tau_{l-1}, \tau_l]$ with $n$
nodes:
\begin{equation*}
  \Omega_n^l = \Bigl\{
    t_j = \tau_{l-1} + (j-1)\omega_n,
    \quad \omega_n = \frac{1}{n-1},
    \quad j = 1, \ldots, n
  \Bigr\}.
\end{equation*}

The Nystr\"{o}m approximation of $T_{lp}$ is:
\begin{equation}
  [T_{n,lp}(\boldsymbol{\Psi}^{(k)}) \phi_p](t) =
  \sum_{j=1}^{n} \omega_n \cdot
  \frac{\partial \kappa_l}{\partial \psi_p}
  \bigl(t, t_j, \psi_p^{(k)}(t_j)\bigr)
  \cdot \phi_p(t_j),
  \label{eq:nystrom}
\end{equation}
and the Nystr\"{o}m error satisfies (from Proposition~2 of
\cite{sedka2025}):
\begin{equation}
  \norm{\DT(\boldsymbol{\Phi}) - \DTn(\boldsymbol{\Phi})}
  \leq \delta_{n,N}
  \xrightarrow{n \to \infty} 0,
  \label{eq:nysterror}
\end{equation}
where
$\delta_{n,N} = \frac{1}{12n^2} \sum_{l=1}^N \sum_{p=1}^N
\sup_{t} \abs{\partial^2 \kappa_l / (\partial \psi_p \partial s)}$.

The TLD linear algebraic system at iteration $k$ is:
\begin{equation}
  \bigl(\lambda I_N - \Hn\bigr) Y^{(k+1)} = b_n^{(k)},
  \label{eq:linsys_tld}
\end{equation}
where $I_N$ is the identity on $\R^{N \cdot n}$, and:
\begin{equation}
  [\Hn]_{lp}(i,j) =
  \omega_n \cdot
  \frac{\partial \kappa_l}{\partial \psi_p}
  \bigl(t_i, t_j, y_{n,p}^{(k)}(j)\bigr).
  \label{eq:Hn}
\end{equation}

\subsubsection{Convergence of TLD}

The following theorem, proved in \cite{sedka2025}, governs
the convergence of TLD.

\begin{theorem}[Convergence of TLD, \cite{sedka2025}]
\label{thm:tld}
Let Assumption~\ref{ass:main} be satisfied, and let
$r = \min(R, 1/(2\gamma\nu))$ where
$\gamma = \max_{1 \leq p \leq N} \gamma_p$.
Then for sufficiently large $n$, there exist
$\alpha_n \in {]0,1[}$ and $\varrho_n > 0$ such that,
if $\boldsymbol{\Psi}_n^{(0)} \in B_{\varrho_n}(\boldsymbol{\Psi})$,
then for all $k \in \N^*$:
\begin{equation}
  \norm{\boldsymbol{\Psi}_n^{(k)} - \boldsymbol{\Psi}}_\mathcal{X}
  \leq \varrho_n
  \left(\frac{1 + \alpha_n}{2}\right)^k
  \xrightarrow{k \to \infty} 0,
  \label{eq:tld_conv}
\end{equation}
where
$\alpha_n = 2\nu\delta_{n,N} / (1 - 2\nu\delta_{n,N})$
and
$\varrho_n = \min\!\bigl(r, (1-\alpha_n)/(2\gamma\nu(1+\alpha_n))\bigr)$.
\end{theorem}

%% -----------------------------------------------------------
\subsection{Physics-Informed Neural Networks}
\label{subsec:pinn}
%% -----------------------------------------------------------

Physics-Informed Neural Networks (PINNs), introduced by
Raissi, Perdikaris, and Karniadakis \cite{raissi2019},
parametrize the unknown function by a neural network and
minimize a loss function built from the residual of the
governing equation evaluated at collocation points.

In the context of integral equations, a PINN approach
defines a network $u_\theta : [0, \tau] \to \R$ with
parameter vector $\theta \in \R^d$, and minimizes:
\begin{equation}
  \mathcal{L}_{\mathrm{PINN}}(\theta) =
  \frac{1}{N_c} \sum_{i=1}^{N_c}
  \left|
    \lambda u_\theta(t_i) -
    \int_0^\tau \kappa(t_i, s, u_\theta(s)) \, ds
    - g(t_i)
  \right|^2,
  \label{eq:pinn_loss}
\end{equation}
where $\{t_i\}_{i=1}^{N_c}$ are collocation points.
While powerful, this approach applied directly to
equation~\eqref{eq:main} with large $\tau$ suffers from
the spectral bias of neural networks \cite{wang2022}, which
causes difficulty in learning complex or oscillatory
solutions over large domains.

In TL-NN, we do not apply PINNs to solve equation
\eqref{eq:main} directly.
Instead, we train a small neural network at each
Newton iteration to approximate a smooth scalar function ---
the kernel derivative $\partial \kappa_l / \partial \psi_p$
at the current iterate --- on a bounded subinterval.
This is a much simpler learning task than the global PINN
approach, and is well within the approximation capacity of
small networks.

%% ============================================================
\section{The TL-NN Algorithm}
\label{sec:tlnn}
%% ============================================================

The TL-NN algorithm modifies the TLD algorithm by replacing
Phase~3 (Nystr\"{o}m discretization) with a neural network
approximation phase.
Phases~1 and 2 are identical to those of \cite{sedka2025}.
We present the complete algorithm here.

\subsection{Phases 1 and 2 (Unchanged)}

Phase 1 applies the domain decomposition
\eqref{eq:decomp}--\eqref{eq:operator}.
Phase 2 defines the Newton--Kantorovich iteration
\eqref{eq:newton}--\eqref{eq:rhs}.
These are exactly as in \cite{sedka2025} and are not modified.

\subsection{Phase 3: Neural Network Approximation}

\subsubsection{The Target Function}

At Newton iteration $k$, for each pair $(l, p)$ with
$l, p \in \{1, \ldots, N\}$, the function to be approximated
is the kernel derivative evaluated at the current iterate:
\begin{equation}
  D_{lp}^{(k)}(t, s) :=
  \frac{\partial \kappa_l}{\partial \psi_p}
  \bigl(t, s, \psi_p^{(k)}(s)\bigr),
  \quad
  (t, s) \in [\tau_{l-1}, \tau_l] \times [\tau_{p-1}, \tau_p].
  \label{eq:target}
\end{equation}
This is precisely the function whose values at quadrature
nodes produce the Nystr\"{o}m matrix $\Hn$ in TLD.
In TL-NN it becomes the training target for the neural
network.

\subsubsection{Network Architecture}

For each pair $(l,p)$ and each Newton iteration $k$, we
define a feedforward neural network:
\begin{equation}
  \ftheta : [\tau_{l-1}, \tau_l] \times [\tau_{p-1}, \tau_p]
  \to \R,
  \label{eq:network}
\end{equation}
with parameter vector $\theta \in \R^d$.
The network has the following architecture:
\begin{itemize}
  \item \textbf{Input layer:} 2 neurons --- the pair $(t, s)$.
  \item \textbf{Hidden layers:} $L = 2$ layers with $m = 20$
    neurons each and activation function $\sigma = \tanh$.
  \item \textbf{Output layer:} 1 neuron --- the scalar
    approximation of $D_{lp}^{(k)}(t, s)$.
\end{itemize}

Explicitly, for $L = 2$ hidden layers:
\begin{equation}
  \ftheta(t, s) =
  W_3 \cdot \sigma\!\left(
    W_2 \cdot \sigma\!\left(
      W_1 \cdot \begin{pmatrix} t \\ s \end{pmatrix} + b_1
    \right) + b_2
  \right) + b_3,
  \label{eq:net_explicit}
\end{equation}
where $W_1 \in \R^{m \times 2}$, $W_2 \in \R^{m \times m}$,
$W_3 \in \R^{1 \times m}$ are weight matrices, and
$b_1, b_2 \in \R^m$, $b_3 \in \R$ are bias vectors.
The parameter vector $\theta$ collects all weights and biases;
for $L = 2$, $m = 20$: $d = 2 \times 20 + 20 \times 20
+ 20 \times 1 + 20 + 20 + 1 = 501$ parameters.

\begin{remark}
  For the smooth kernels in our numerical examples, $L = 2$
  and $m = 20$ are sufficient. For more complex or oscillatory
  kernels, the architecture can be enlarged without modifying
  the algorithm structure.
\end{remark}

\subsubsection{Training Data}

The training dataset for $\ftheta$ consists of input-output
pairs evaluated at the grid nodes:
\begin{equation}
  \mathcal{T}^{(k,lp)} = \left\{
    \bigl( (t_i, t_j),\; D_{lp}^{(k)}(t_i, t_j) \bigr)
    : t_i \in \Omega_n^l,\; t_j \in \Omega_n^p
  \right\},
  \label{eq:training_data}
\end{equation}
a dataset of $n^2$ pairs.
The target values $D_{lp}^{(k)}(t_i, t_j)$ are computed by
direct evaluation of the kernel derivative at the current
iterate $\psi_p^{(k)}$ --- the same values that the TLD code
already computes to build $\Hn$.
No additional function evaluations are required.

\subsubsection{Loss Function and Training}

The network $\ftheta$ is trained by minimizing the mean
squared error loss:
\begin{equation}
  \mathcal{L}(\theta) =
  \frac{1}{n^2} \sum_{i=1}^{n} \sum_{j=1}^{n}
  \Bigl[
    \ftheta(t_i, t_j) - D_{lp}^{(k)}(t_i, t_j)
  \Bigr]^2.
  \label{eq:loss}
\end{equation}

Training is performed using the Adam optimizer
\cite{kingma2015} with learning rate $\eta = 10^{-3}$,
running for at most $N_{\max}$ epochs.
The parameter update at each epoch is:
\begin{equation}
  \theta \leftarrow \theta - \eta \cdot \nabla_\theta
  \mathcal{L}(\theta),
  \label{eq:adam}
\end{equation}
with Adam moment estimates applied as described in
\cite{kingma2015}.
Training stops when
$\norm{\nabla_\theta \mathcal{L}(\theta)} < \varepsilon_{\rm train}$
or when $N_{\max}$ epochs are reached.

\begin{remark}
  Since $D_{lp}^{(k)}$ is a smooth scalar function of two
  variables on a bounded domain, and since the network is
  small (501 parameters), training to high accuracy typically
  requires between 500 and 3000 epochs for the examples
  in Section~\ref{sec:numerics}.
\end{remark}

\subsubsection{Neural Network Approximation Error}

After training, the approximation quality of $\ftheta$ is
measured by:
\begin{equation}
  \epsNN^{(k,lp)} = \sup_{(t,s) \in
  [\tau_{l-1}, \tau_l] \times [\tau_{p-1}, \tau_p]}
  \left|
    \ftheta(t, s) - D_{lp}^{(k)}(t, s)
  \right|.
  \label{eq:nn_error}
\end{equation}

The global neural network approximation error at iteration
$k$ is:
\begin{equation}
  \epsNN = \max_{1 \leq l, p \leq N} \epsNN^{(k,lp)}.
  \label{eq:nn_error_global}
\end{equation}
This quantity plays the same role in the convergence analysis
of TL-NN as $\delta_{n,N}$ plays in TLD.

\subsubsection{Building the TL-NN Matrix}

The matrix $\HNN$ is built from neural network evaluations:
\begin{equation}
  [\HNN]_{lp}(i,j) =
  \omega_n \cdot \ftheta(t_i, t_j),
  \quad l, p = 1, \ldots, N,
  \quad i, j = 1, \ldots, n.
  \label{eq:Hnn}
\end{equation}

Comparing \eqref{eq:Hnn} with \eqref{eq:Hn}, we see that
$\HNN$ and $\Hn$ have identical dimensions, identical block
structure, and identical quadrature weights $\omega_n$.
The only difference is that the kernel derivative values
$D_{lp}^{(k)}(t_i, t_j)$ are produced by neural network
evaluation in TL-NN rather than direct computation in TLD.

\subsection{The TL-NN Linear System}

\textbf{Step 1.} Solve the linear algebraic system:
\begin{equation}
  \bigl(\lambda I_N - \HNN\bigr) Y^{(k+1)} = b_{NN}^{(k)},
  \label{eq:linsys_tlnn}
\end{equation}
where the right-hand side is:
\begin{equation}
  b_{NN,l}^{(k)}(i) =
  - \sum_{p=1}^{N} \sum_{j=1}^{n}
    \omega_n \cdot \ftheta(t_i, t_j)
    \cdot y_{n,p}^{(k)}(j)
  + \sum_{p=1}^{N} \sum_{z=1}^{m}
    \omega_m \cdot \kappa_l\!\bigl(t_i, t_z,
    \psi_{n,p}^{(k)}(t_z)\bigr)
  + g_l(t_i).
  \label{eq:rhs_tlnn}
\end{equation}

\textbf{Step 2.} Recover the iterate solution by the natural
interpolation formula:
\begin{equation}
  \lambda \psi_{n,l}^{(k+1)}(t) =
  \sum_{p=1}^{N} \sum_{j=1}^{n}
  \omega_n \cdot \ftheta(t, t_j)
  \cdot \bigl(y_{n,p}^{(k+1)}(j) - y_{n,p}^{(k)}(j)\bigr)
  + \sum_{p=1}^{N} \sum_{z=1}^{m}
  \omega_m \cdot \kappa_l\!\bigl(t, t_z,
  \psi_{n,p}^{(k)}(t_z)\bigr)
  + g_l(t).
  \label{eq:interp_tlnn}
\end{equation}

\subsection{Stopping Condition and Error}

The Newton iteration stops when:
\begin{equation}
  E_n^{(k)} = \sum_{l=1}^{N} \max_{1 \leq i \leq n}
  \left|
    \psi_{n,l}^{(k+1)}(t_i) - \psi_{n,l}^{(k)}(t_i)
  \right| \leq 10^{-9}.
  \label{eq:stop}
\end{equation}

The approximation error against the exact solution
$\psi^{\rm ext}$ is:
\begin{equation}
  E_{n,N} = \sum_{l=1}^{N} \max_{1 \leq i \leq n}
  \left|
    \psi_l^{\rm ext}(t_i) - \psi_{n,l}^{(k+1)}(t_i)
  \right|.
  \label{eq:error}
\end{equation}

\subsection{Algorithm Summary}

The complete TL-NN algorithm is presented in
Algorithm~\ref{alg:tlnn}.

\begin{algorithm}[H]
\caption{TL-NN Algorithm}
\label{alg:tlnn}
\begin{algorithmic}[1]
\Require $n, N, m, \kappa, g, \lambda, \tau$,
  $N_{\max}$, $\varepsilon_{\rm train}$,
  $\psi^{\rm ext}$ (for error computation)
\Ensure $E_n^{(k)}$, $E_{n,N}$,
  $\boldsymbol{\Psi}_n^{(k+1)}$
\State \textbf{Phase 1 (Transform):}
  Compute subintervals $[\tau_{p-1}, \tau_p]$
  and grids $\Omega_n^l$, $\Omega_m^l$
\State Initialize $\boldsymbol{\Psi}_n^{(0)} \leftarrow \mathbf{0}$,
  $k \leftarrow 0$, $E_n^{(0)} \leftarrow 1$
\While{$E_n^{(k)} > 10^{-9}$}
  \State \textbf{Phase 2 (Linearize):}
    Compute $F(\boldsymbol{\Psi}_n^{(k)})$ and
    evaluate $D_{lp}^{(k)}(t_i, t_j)$ for all $l, p, i, j$
  \State \textbf{Phase 3 (Neural Network):}
  \For{each pair $(l, p)$, $l, p = 1, \ldots, N$}
    \State Build training set
      $\mathcal{T}^{(k,lp)}$ from \eqref{eq:training_data}
    \State Initialize network $\ftheta$ with random $\theta$
    \For{epoch $= 1$ to $N_{\max}$}
      \State Compute loss $\mathcal{L}(\theta)$
        from \eqref{eq:loss}
      \State Update $\theta$ with Adam optimizer
      \If{$\norm{\nabla_\theta \mathcal{L}} < \varepsilon_{\rm train}$}
        \State \textbf{break}
      \EndIf
    \EndFor
    \State Build block $[\HNN]_{lp}$
      from \eqref{eq:Hnn}
  \EndFor
  \State Build right-hand side $b_{NN}^{(k)}$
    from \eqref{eq:rhs_tlnn}
  \State Solve $(\lambda I_N - \HNN) Y^{(k+1)}
    = b_{NN}^{(k)}$
  \State Recover $\boldsymbol{\Psi}_n^{(k+1)}$
    from \eqref{eq:interp_tlnn}
  \State Compute $E_n^{(k)}$ from \eqref{eq:stop}
  \State $k \leftarrow k + 1$
\EndWhile
\State Compute $E_{n,N}$ from \eqref{eq:error}
\end{algorithmic}
\end{algorithm}

%% ============================================================
\section{Convergence Discussion}
\label{sec:convergence}
%% ============================================================

The convergence framework of TLD carries over directly to
TL-NN. The key structural observation is the following.

\begin{proposition}[Analog of Proposition~4 in \cite{sedka2025}]
\label{prop:invertible}
Let Assumption~\ref{ass:main} hold and let
$r = \min(R, 1/(2\gamma\nu))$.
Suppose the neural network approximation error satisfies
$\epsNN < 1/(2\nu)$.
Then for all $\boldsymbol{\Phi} \in B_r(\boldsymbol{\Psi})$,
the operator $(\lambda I - \DTNN(\boldsymbol{\Phi}))$ is
invertible and:
\begin{equation}
  \sup_{\boldsymbol{\Phi} \in B_r(\boldsymbol{\Psi})}
  \norm{(\lambda I - \DTNN(\boldsymbol{\Phi}))^{-1}}
  \leq 2\nu(1 + \alphNN),
  \label{eq:inv_bound}
\end{equation}
where:
\begin{equation}
  \alphNN = \frac{2\nu\epsNN}{1 - 2\nu\epsNN}.
  \label{eq:alphNN}
\end{equation}
\end{proposition}

\begin{remark}
  Proposition~\ref{prop:invertible} follows from the same
  geometric series argument as Proposition~4 in
  \cite{sedka2025}, with $\delta_{n,N}$ replaced by
  $\epsNN$.
  The proof is omitted as it is formally identical.
\end{remark}

\begin{theorem}[Convergence of TL-NN]
\label{thm:tlnn}
Under the conditions of Proposition~\ref{prop:invertible},
there exist $\alphNN \in {]0,1[}$ and $\varrhoNN > 0$ such
that, if $\boldsymbol{\Psi}_{NN}^{(0)} \in
B_{\varrhoNN}(\boldsymbol{\Psi})$, then for all $k \in \N^*$:
\begin{equation}
  \norm{\boldsymbol{\Psi}_{NN}^{(k)} -
  \boldsymbol{\Psi}}_\mathcal{X}
  \leq \varrhoNN
  \left(\frac{1 + \alphNN}{2}\right)^k
  \xrightarrow{k \to \infty} 0,
  \label{eq:tlnn_conv}
\end{equation}
where $\alphNN$ is given by \eqref{eq:alphNN} and:
\begin{equation}
  \varrhoNN = \min\!\left(r,\;
  \frac{1 - \alphNN}{2\gamma\nu(1 + \alphNN)}\right).
  \label{eq:varrhoNN}
\end{equation}
\end{theorem}

\begin{remark}[Comparison with TLD]
  Comparing Theorem~\ref{thm:tld} and Theorem~\ref{thm:tlnn},
  the convergence rates of TLD and TL-NN are governed by
  $(1+\alpha_n)/2$ and $(1+\alphNN)/2$ respectively.
  Since both $\alpha_n$ and $\alphNN$ are increasing
  functions of their respective errors, TL-NN converges
  faster than TLD whenever $\epsNN < \delta_{n,N}$.
  
  For Nystr\"{o}m on a uniform grid, $\delta_{n,N}$ decreases
  as $O(1/n^2)$ and depends on the second derivative of the
  kernel with respect to $s$.
  For oscillatory or rapidly varying kernels on large
  intervals, this derivative is large and $n$ must be
  correspondingly large to make $\delta_{n,N}$ small.
  A neural network trained on well-chosen collocation points
  adapts its approximation capacity to the local behavior of
  $D_{lp}^{(k)}$, potentially achieving smaller $\epsNN$
  with fewer evaluation points.
  The numerical experiments in Section~\ref{sec:numerics}
  assess this comparison quantitatively.
\end{remark}

%% ============================================================
\section{Numerical Results}
\label{sec:numerics}
%% ============================================================

We apply TL-NN to the three benchmark examples from
\cite{sedka2025}, using identical parameters to allow
direct comparison.
For all experiments: $N = \tau/10$ subintervals (so that
each subinterval has length 10), $m = 10n$ fine quadrature
nodes for the integral evaluation, $\lambda = 10$,
$N_{\max} = 3000$ training epochs, and
$\varepsilon_{\rm train} = 10^{-8}$.
The neural network architecture is $L=2$, $m=20$, $\sigma
= \tanh$, trained with Adam at learning rate $\eta = 10^{-3}$.

All computations are performed in Python using PyTorch for
the neural network component and NumPy for the linear algebra.
The TLD results are reproduced from \cite{sedka2025} for
reference.

\subsection{Example 1}

Consider:
\begin{equation}
  \lambda\psi(t) = \int_0^\tau
  \frac{(t+s)^3 \psi^2(s)}{(4 + \sin s - 2\sin^2(2s))^2}
  \, ds + g(t),
  \quad t \in [0, \tau],
  \label{eq:ex1}
\end{equation}
with exact solution $\psi_{\rm ext}(t) = 3 + \sin t + \cos(4t)$
and $g(t) = \lambda \psi_{\rm ext}(t) - \frac{1}{4}
[(t+\tau)^4 - t^4]$.

\begin{table}[H]
\centering
\caption{Errors $E_{n,N}$ for Example~1 ($m = 10n$,
$\lambda = 10$): TLD \cite{sedka2025} vs TL-NN (this work).}
\label{tab:ex1}
\begin{tabular}{c cc cc cc}
\toprule
& \multicolumn{2}{c}{$\tau = 20$}
& \multicolumn{2}{c}{$\tau = 50$}
& \multicolumn{2}{c}{$\tau = 80$} \\
\cmidrule(lr){2-3}\cmidrule(lr){4-5}\cmidrule(lr){6-7}
$n$ & TLD & TL-NN & TLD & TL-NN & TLD & TL-NN \\
\midrule
5   & \textit{[TLD]} & \textit{[TL-NN]} &
      \textit{[TLD]} & \textit{[TL-NN]} &
      \textit{[TLD]} & \textit{[TL-NN]} \\
10  & & & & & & \\
30  & & & & & & \\
50  & & & & & & \\
100 & & & & & & \\
\bottomrule
\end{tabular}
\end{table}

\begin{remark}
  Table~\ref{tab:ex1} will be completed with numerical
  results upon implementation.
  The TLD values are available from Table~1 of
  \cite{sedka2025}.
\end{remark}

\subsection{Example 2}

Consider:
\begin{equation}
  \lambda\psi(t) = \int_0^\tau
  \cos(3s) \, e^{(t+3s)/30} \psi^3(s)
  \, ds + g(t),
  \quad t \in [0, \tau],
  \label{eq:ex2}
\end{equation}
with exact solution $\psi_{\rm ext}(t) = e^{-t/30}\sin(3t)$
and $g(t) = \lambda e^{-t/30}\sin(3t)
- \frac{e^{t/30}}{16}\sin^4(3t)$.

\begin{table}[H]
\centering
\caption{Errors $E_{n,N}$ for Example~2 ($m = 10n$):
TLD \cite{sedka2025} vs TL-NN (this work).}
\label{tab:ex2}
\begin{tabular}{c cc cc cc}
\toprule
& \multicolumn{2}{c}{$\tau = 10$}
& \multicolumn{2}{c}{$\tau = 50$}
& \multicolumn{2}{c}{$\tau = 80$} \\
\cmidrule(lr){2-3}\cmidrule(lr){4-5}\cmidrule(lr){6-7}
$n$ & TLD & TL-NN & TLD & TL-NN & TLD & TL-NN \\
\midrule
5   & & & & & & \\
10  & & & & & & \\
30  & & & & & & \\
50  & & & & & & \\
100 & & & & & & \\
\bottomrule
\end{tabular}
\end{table}

\subsection{Example 3}

Consider:
\begin{equation}
  \lambda\psi(t) = \frac{1}{4}\int_0^\tau
  t e^{-ts}\bigl[2(1-\cos(\pi s)) + \psi^2(s)\bigr]
  \, ds + g(t),
  \quad t \in [0, \tau],
  \label{eq:ex3}
\end{equation}
with exact solution $\psi_{\rm ext}(t) = 2\cos(\pi t/2)$
and $g(t) = 2\lambda\cos(\pi t/2) - (1 - e^{-t\tau})$.

\begin{table}[H]
\centering
\caption{Errors $E_{n,N}$ for Example~3 ($m = 10n$,
$\lambda = 10$): TLD \cite{sedka2025} vs TL-NN (this work).}
\label{tab:ex3}
\begin{tabular}{c cc cc cc}
\toprule
& \multicolumn{2}{c}{$\tau = 20$}
& \multicolumn{2}{c}{$\tau = 50$}
& \multicolumn{2}{c}{$\tau = 100$} \\
\cmidrule(lr){2-3}\cmidrule(lr){4-5}\cmidrule(lr){6-7}
$n$ & TLD & TL-NN & TLD & TL-NN & TLD & TL-NN \\
\midrule
5   & & & & & & \\
10  & & & & & & \\
30  & & & & & & \\
50  & & & & & & \\
80  & & & & & & \\
\bottomrule
\end{tabular}
\end{table}

\subsection{Discussion}

\textcolor{red}{[TO BE COMPLETED AFTER NUMERICAL EXPERIMENTS]}

This section will contain:
\begin{itemize}
  \item A quantitative comparison of $E_{n,N}$ between
    TLD and TL-NN across all examples and all values of
    $\tau$ and $n$.
  \item A discussion of the training cost of the neural
    network component versus the quadrature cost of Nystr\"{o}m.
  \item An analysis of when TL-NN outperforms TLD and when
    TLD remains preferable, with reference to the convergence
    rate comparison in Remark~\ref{thm:tlnn}.
  \item Figures showing approximate solutions versus exact
    solutions for selected values of $n$ and $\tau$.
  \item The variation of $\log_{10}(E_n^{(k)})$ over
    Newton iterations, confirming linear convergence rate.
\end{itemize}

%% ============================================================
\section{Conclusion}
\label{sec:conclusion}
%% ============================================================

We have introduced TL-NN, a hybrid numerical method for
nonlinear Fredholm integral equations on large intervals,
obtained by replacing the Nystr\"{o}m discretization phase
of the TLD algorithm \cite{sedka2025} with a neural network
approximator trained at each Newton--Kantorovich iteration.

The theoretical framework of TLD is preserved in its entirety.
The convergence rate of TL-NN is governed by the same formula
as TLD, with the Nystr\"{o}m approximation error $\delta_{n,N}$
replaced by the neural network approximation error $\epsNN$.
On large intervals with complex or oscillatory kernels, neural
networks can achieve smaller $\epsNN$ than the corresponding
$\delta_{n,N}$ for a fixed number of evaluation points,
because they distribute their approximation capacity
adaptively rather than uniformly.

Numerical experiments on three benchmark examples from
\cite{sedka2025} demonstrate that TL-NN produces results
comparable to TLD in terms of accuracy, with characteristics
that make it particularly competitive for large $\tau$.

\textbf{Future work.}
Several extensions of TL-NN are natural directions for
subsequent research:
\begin{enumerate}
  \item \textbf{Weakly singular kernels.} The extension of
    TL-NN to kernels with weak singularities, as studied
    in \cite{sedka2022b} for TLD, represents a natural
    and important generalization. Neural networks are known
    to handle certain types of singular behavior better
    than uniform quadrature grids.
  \item \textbf{Systems of equations.} The extension to
    systems of nonlinear Fredholm integral equations,
    as treated in \cite{sedka2022a} and \cite{sedka2022b},
    follows directly from the present formulation.
  \item \textbf{Convergence theorem for the neural network
    training.} A rigorous bound on $\epsNN$ in terms of
    the neural network architecture, the number of training
    epochs, and the properties of the kernel $\kappa$ would
    place TL-NN on the same fully rigorous theoretical
    footing as TLD. This constitutes the main open
    theoretical question left by the present work.
  \item \textbf{TLD-XPINN.} A deeper integration of the
    TLD domain decomposition with Extended Physics-Informed
    Neural Networks \cite{jagtap2020}, assigning a separate
    network to each subinterval and training them
    simultaneously through the global Fredholm coupling,
    represents a more ambitious hybrid that we plan to
    investigate in subsequent work.
\end{enumerate}

%% ============================================================
%% ACKNOWLEDGEMENTS
%% ============================================================
\section*{Acknowledgements}

The authors thank the General Direction of Scientific
Research and Technological Development (DGRSDT) of the
Ministry of Higher Education and Scientific Research of
Algeria for supporting this research.

%% ============================================================
%% BIBLIOGRAPHY
%% ============================================================
\bibliographystyle{elsarticle-num}

\begin{thebibliography}{99}

\bibitem{sedka2025}
I. Sedka, A. Khellaf,
Toward efficient numerical solutions of non-linear integral
equations with TLD algorithm,
\textit{Izvestiya: Mathematics} 89(2) (2025) 341--358.
\url{https://doi.org/10.4213/im9613e}

\bibitem{sedka2023}
I. Sedka, A. Khellaf, M.Z. Aissaoui,
New algorithm to solve nonlinear functional equations
applying linearization then double discretization scheme
(L.D.D),
\textit{Kuwait Journal of Science} 50(2) (2023) 65--74.
\url{https://doi.org/10.1016/j.kjs.2023.02.010}

\bibitem{sedka2022a}
I. Sedka, S. Lemita, M.Z. Aissaoui,
Linearization-Discretization process to solve systems of
nonlinear Fredholm integral equations in an
infinite-dimensional context,
\textit{Advances in the Theory of Nonlinear Analysis and
its Applications} 6(4) (2022) 547--564.
\url{https://doi.org/10.31197/atnaa.998275}

\bibitem{sedka2022b}
I. Sedka, A. Khellaf, S. Lemita, M.Z. Aissaoui,
New Algorithm on Linearization-Discretization Solving
Systems of Nonlinear Integro-Differential Fredholm
Equations,
\textit{Boletim da Sociedade Paranaense de Matem\'{a}tica}
42 (2024) 1--17.
\url{https://doi.org/10.5269/bspm.63480}

\bibitem{lemita2020}
S. Lemita, H. Guebbai, I. Sedka, M.Z. Aissaoui,
New method for the numerical solution of the Fredholm
linear integral equation on a large interval,
\textit{Russian Universities Reports. Mathematics}
25(132) (2020) 387--400.

\bibitem{atkinson1997}
K.E. Atkinson,
\textit{The Numerical Solution of Integral Equations of
the Second Kind},
Cambridge University Press, Cambridge, 1997.

\bibitem{grammont2013}
L. Grammont,
Nonlinear integral equation of the second kind:
a new version of Nystr\"{o}m method,
\textit{Numerical Functional Analysis and Optimization}
34(5) (2013) 496--515.

\bibitem{grammont2014}
L. Grammont, M. Ahues, F.D. D'Almeida,
For nonlinear infinite dimensional equations, which to
begin with: linearization or discretization?
\textit{Journal of Integral Equations and Applications}
26(3) (2014) 413--436.

\bibitem{khellaf2022}
A. Khellaf, M.Z. Aissaoui,
New theoretical conditions for solving functional
nonlinear equations by linearization then discretization,
\textit{International Journal of Nonlinear Analysis and
Applications} 13(1) (2022) 2857--2869.

\bibitem{raissi2019}
M. Raissi, P. Perdikaris, G.E. Karniadakis,
Physics-informed neural networks: A deep learning
framework for solving forward and inverse problems
involving nonlinear partial differential equations,
\textit{Journal of Computational Physics} 378 (2019)
686--707.
\url{https://doi.org/10.1016/j.jcp.2018.10.045}

\bibitem{lu2021}
L. Lu, P. Jin, G. Pang, Z. Zhang, G.E. Karniadakis,
Learning nonlinear operators via DeepONet based on the
universal approximation theorem of operators,
\textit{Nature Machine Intelligence} 3 (2021) 218--229.
\url{https://doi.org/10.1038/s42256-021-00302-5}

\bibitem{kovachki2023}
N. Kovachki, Z. Li, B. Liu, K. Bhattacharya,
A. Stuart, A. Anandkumar,
Neural operator: Learning maps between function spaces
with applications to PDEs,
\textit{Journal of Machine Learning Research}
24(89) (2023) 1--97.

\bibitem{kingma2015}
D.P. Kingma, J.L. Ba,
Adam: A method for stochastic optimization,
in: \textit{Proceedings of the 3rd International
Conference on Learning Representations (ICLR 2015)},
arXiv:1412.6980, 2015.

\bibitem{wang2022}
S. Wang, X. Yu, P. Perdikaris,
When and why PINNs fail to train: A neural tangent
kernel perspective,
\textit{Journal of Computational Physics} 449 (2022)
110768.

\bibitem{jagtap2020}
A.D. Jagtap, E. Kharazmi, G.E. Karniadakis,
Conservative physics-informed neural networks on discrete
domains for conservation laws: Applications to forward
and inverse problems,
\textit{Computer Methods in Applied Mechanics and
Engineering} 365 (2020) 113028.

\bibitem{guan2022}
Y. Guan, T. Fang, D. Zhang, C. Jin,
Solving Fredholm integral equations using deep learning,
\textit{International Journal of Applied and Computational
Mathematics} 8(2) (2022) 87.

\bibitem{ladopoulos2000}
E.G. Ladopoulos,
\textit{Singular Integral Equations: Linear and
Non-Linear Theory and its Applications in Science and
Engineering},
Springer, Berlin, 2000.

\bibitem{fernandez2017}
J.A.E. Fern\'{a}ndez, M. Ver\'{o}n,
\textit{Newton's Method: An Updated Approach of
Kantorovich's Theory},
Birkh\"{a}user/Springer, Cham, 2017.

\bibitem{zabrejko1996}
P.P. Zabrejko,
The mean value theorem for differentiable mappings in
Banach spaces,
\textit{Integral Transforms and Special Functions}
4(1--2) (1996) 153--162.

\end{thebibliography}

\end{document}
